\section{Background}
\label{sec:background}

The problem of accountability in autonomous systems is not new, but its urgency has escalated sharply with the deployment of multi-agent AI architectures in high-stakes environments. This section reviews four foundational threads that converge in the design of the Bailout Protocol: the principle of separation of concerns as a basis for architectural accountability; accountability as a first-class property in multi-agent AI systems; human-in-the-loop design as the dominant paradigm for maintaining human oversight; and fail-safe design as established in safety-critical systems engineering. We then situate the BX3 Framework's three-layer model as an architecture that synthesizes these threads into a unified accountability design.

\subsection{Separation of Concerns and Accountability Architecture}

The conceptual foundation for treating accountability as an architectural property rather than a governance afterthought traces to Dijkstra's articulation of the \textit{separation of concerns} principle. In \cite{dijkstra1982}, Dijkstra argued that sound system design requires deliberate boundaries between functional domains, such that each concern can be reasoned about and enforced independently. When concerns are entangled, failures in one domain propagate unpredictably into others, making responsibility unassignable and remediation impossible. Dijkstra's insight, originally applied to software module boundaries, extends naturally to the boundary between autonomous machine action and human accountability. If a machine actor can modify the conditions under which it is evaluated, or if the boundary between machine and human authority is fluid, accountability dissolves. The separation of concerns demands that the authority to set objectives, the authority to propose actions, and the authority to execute physical operations occupy distinct, immutable layers of the system architecture. The BX3 Framework takes this principle as its foundational commitment: each layer has a fixed role, and no layer can unilaterally modify the constraints of another.

\subsection{Accountability in Multi-Agent AI Systems}

Classical multi-agent systems (MAS) research established that when multiple autonomous actors operate concurrently, attributing responsibility for collective outcomes requires explicit architectural support. As surveyed in \cite{trist1981}, the challenge is fundamental: agents are defined by their autonomy --- the capacity to act without direct human instruction --- which is precisely the property that severs the human--action chain needed for accountability. In traditional MAS, accountability has been addressed through role assignment (agents play defined roles with associated duties), through voting or consensus protocols (collective decisions distribute or obscure individual responsibility), and through organizational norms (informal or formal conventions about who is responsible for what). Each of these approaches is susceptible to the \textit{accountability gap}: the scenario in which an autonomous action produces a harmful outcome but no human actor can reconstruct, explain, or be held responsible for the decision chain that led there.

Modern AI agent frameworks have intensified this problem. Large language model--based agents operate probabilistically rather than deterministically, making the input--output relationship non-invertible and the decision logic opaque even to their designers. When multiple such agents are composed in orchestration pipelines, the accountability chain fractures across model boundaries, prompt contexts, and tool invocations that no single human reviews before execution. Recent regulatory frameworks, including the EU AI Act and the NIST AI Risk Management Framework, now require that high-risk AI systems maintain ``human oversight'' and ``decision provenance'' not as best practices but as mandatory compliance obligations \cite{trist1981}. These mandates presuppose an architecture that can satisfy them --- a structure in which every autonomous action either originates from or can be traced to a human accountability anchor. The Bailout Protocol is designed as that architectural substrate.

\subsection{Human-in-the-Loop Design}

Human-in-the-loop (HITL) is the dominant paradigm for incorporating human oversight into autonomous systems. In HITL design, a human operator sits at one or more points in the control loop, reviewing, approving, or correcting agent actions before or after execution. The paradigm exists on a spectrum: at one end, humans approve every action in a strict workflow gate; at the other, humans are notified post-hoc of actions already taken, retaining only the ability to investigate and remediate rather than prevent.

The fundamental tension in HITL design is between the efficiency of automation and the fidelity of human oversight. Pure automation achieves speed and scale but sacrifices accountability; full human review preserves accountability but negates the autonomy that makes multi-agent systems valuable. Existing HITL frameworks address this tension through confidence thresholds --- autonomous action proceeds when the system's confidence exceeds a configurable level, and human review is triggered only below that threshold. However, this approach has a structural flaw in the context of LLM-based agents: the confidence score produced by a model is not a reliable indicator of correctness with respect to the user's actual objectives, because the model's training objective (next-token prediction) is decoupled from the operational objective (correct action in context). A model can express high confidence while being wrong about what the user actually intended.

The BX3 Framework inverts this logic. Rather than triggering human review based on a machine-generated confidence signal, the Bailout Protocol treats human review as the \textit{default state} for any unresolvable condition, and uses machine confidence only as a signal internal to the machine layers. If the Bounds Engine cannot produce a proposal that satisfies all constraints within its authority, it does not guess --- it escalates unconditionally. This architectural commitment eliminates the failure mode in which a highly confident but incorrect agent proceeds to harmful action because no threshold was crossed.

HITL design also confronts the \textit{abstraction leakage} problem \cite{trist1981}: when a human operator is inserted into a running system to investigate or override an agent's action, the operator often loses the contextual information needed to make an informed decision. The hierarchical context of the agent's reasoning, the state of the broader multi-agent system, and the chain of prior decisions that led to the current state are distributed across system logs that no single interface presents coherently. The BX3 Framework addresses this through the Forensic Ledger (Section~\ref{sec:forensic-ledger}), which maintains synchronized, plane-separated records of every decision state, ensuring that a human receiving a Bailout signal has complete diagnostic context available for decision-making.

\subsection{Fail-Safe Design in Safety-Critical Systems}

The engineering of fail-safe behavior in safety-critical systems provides the most mature body of knowledge for designing systems that behave predictably under unexpected conditions. A fail-safe system is one that transitions to a defined safe state upon failure, as opposed to a fail-dangerous system that continues operating in an unsafe manner. The canonical example is a nuclear reactor scram: upon detecting anomalous conditions, the control system inserts the control rods regardless of any other operational considerations, because the cost of continued operation under uncertainty exceeds the cost of a controlled shutdown.

In autonomous vehicles and cyber-physical systems, fail-safe design has been extensively studied. The Simplex architecture \cite{bansal2019} introduced the concept of a \textit{safety controller} that runs in parallel with a high-performance primary controller and can commandeer system execution when the primary controller's proposed action would violate a formally verified safety constraint. This decoupled architecture --- performance through the primary controller, safety through an independent, verifiable safety controller --- is the direct progenitor of the BX3 three-layer model. In the Simplex model, the safety controller is a software component; in BX3, the analogous function is performed by the Fact Layer, which enforces physical and operational constraints that no Bounds Engine can override, and by the Bailout Protocol, which ensures that unresolved states reach a human safety anchor rather than being resolved by further machine inference.

More recent work \cite{bansal2019} extends the fail-safe principle to learning-enabled systems, where the primary controller may be a neural network with unpredictable failure modes. The \textit{Perception Simplex} architecture introduces a fault-tolerant safety layer that can detect when the perception system has entered a state of sensor or model failure and override the mission controller with a safe fallback trajectory. The architectural pattern is consistent: a dedicated, verifiable safety layer that can observe and interrupt the primary operational controller, with provable guarantees about the conditions under which override occurs. BX3's Fact Layer encodes this pattern by enforcing that no physical actuator receives a command that has not been validated against a deterministic model of safe operation. The Sandbox Gate further extends this pattern by requiring digital-twin simulation of any proposed intervention before physical execution is permitted.

A critical distinction separates BX3's fail-safe design from prior work. In the Simplex family of architectures, the safety controller is itself a machine component --- albeit a simpler, more formally verifiable one. The override decision is therefore still a machine decision, subject to the same uncertainty and failure modes as the primary controller, though reduced in complexity. In the BX3 Framework, the Bailout Protocol terminates at a human decision-maker, not a machine fallback. The human is not a final-tier machine in the escalation hierarchy; the human is an accountability anchor who receives complete diagnostic context and makes a genuine, uncomputable judgment call. This reflects the principle articulated in \cite{wiener1948}: that machines operating in open environments with genuine uncertainty must ultimately defer to human judgment not because humans are computationally superior, but because accountability for consequential decisions in the real world is a social and legal institution that cannot be instantiated in software.

\subsection{The BX3 Framework: A Three-Layer Accountability Architecture}

The BX3 Framework synthesizes the preceding threads into a unified architectural model with three immutable functional layers and two enforcement mechanisms. The layers are:

\begin{enumerate}

\item \textbf{Purpose Layer (Human Accountability Anchor):} The topmost layer consists exclusively of human actors with defined organizational authority. Each node in the system has a Purpose --- a specification of the objective it serves --- that is authored and owned by a named human. The Purpose Layer never proposes actions autonomously; it authorizes, overrides, or escalates. Every action traceable to the system must have a Purpose Layer owner.

\item \textbf{Bounds Engine Layer (Constraint Enforcement):} The middle layer contains autonomous actors that propose actions within the authority granted by the Purpose Layer. The Bounds Engine's role is to evaluate proposed actions against the constraints encoded in its operational mandate and to either approve the action (if all constraints are satisfied), reject the action (if constraints are violated and no alternative satisfies them), or invoke the Bailout Protocol (if the situation exceeds its authority to resolve). The Bounds Engine is explicitly a proposal engine, not an execution engine --- it never drives physical actuators directly.

\item \textbf{Fact Layer (Deterministic Ground):} The bottom layer encodes the deterministic, verifiable ground truth of the physical or digital system under control. The Fact Layer is the only layer with direct access to physical actuators and data sources. It enforces that all commands reaching actuators conform to the constraints the Fact Layer's deterministic model defines as safe. The Fact Layer is designed to produce identical outputs for identical inputs, enabling formal verification and forensic reconstruction.

\end{enumerate}

The two enforcement mechanisms are the \textbf{Sandbox Gate} and the \textbf{Bailout Protocol}. The Sandbox Gate requires that any Bounds Engine proposal with physical-world consequences be simulated in a digital twin environment and reviewed by the appropriate Fact Layer gate before physical execution is unlocked. The Bailout Protocol is the subject of this paper: it guarantees that any state the Bounds Engine cannot resolve within its current authority escalates directly to the Purpose Layer human anchor, bypassing all intermediate machine actors, with complete diagnostic context. The system fails upward into human consciousness --- it never fails downward into algorithmic chaos.

This architecture satisfies the accountability requirements established by the NIST AI Risk Management Framework and the EU AI Act not through governance documentation, but through architectural enforcement: every autonomous action either originates from a Purpose Layer authorization or triggers a Bailout that terminates at one. The separation of concerns is not a convention; it is a physical constraint enforced by the layer architecture. The human accountability anchor is not a last resort; it is the required terminal state of every unresolvable exception.
